“Jager wordt mijn visser!”,
zei hij als opfrisser.

“Ik vroeg je om mensen!
Respecteer de wensen

die ik je eerder vroeg.
Zoals Jezus opdroeg

mensen te gaan vissen.
Die zich vergewissen

van mijn heilige leer.”
De profeet wilde meer

aanhangers uit het net
dat sluw was uitgezet.

De jager was het zat
en zei: “U kunt me wat!

Ik jaag met veel gevaar
bent U dat wel gewaar?

Wist u dat zelfs zwanen,
die zich bedreigd wanen,

hard met hun vleugels slaan?
Je dan beter kunt gaan?”

De herder verzuchtte
toen hij zijn hart luchtte:

“Een volger van de Heer,
mekkert niet elke keer

bij een kleine drempel,
maar zet ferm zijn stempel

en zal zeker geen zwaan
vervaard uit de weg gaan.

Dus mocht men me missen
dan ben ik gaan vissen!

Want met dit spraakgebruik
vul ik geen lege buik

maar voed ik slechts jouw angst!”
Voor een goede visvangst

wees men op de loopbrug.
De profeet kwam terug

van top tot teen kletsnat.
Van voorhoofd tot kontgat

bedekt met eendenkroos.
Hij liep dwalend en broos

met tussen zijn kleren
een paar witte veren.

De jager als jachtvisser

Het neologisme "jachtvisser" kan op verschillende manieren geïnterpreteerd worden, afhankelijk van de context. Hier zijn enkele mogelijke betekenissen:

1. **Hybride levensstijl of beroep**
Een "jachtvisser" zou kunnen verwijzen naar iemand die een combinatie van jagen en vissen beoefent als een bedrijfsmatig beroep of hobby. Deze persoon leeft als het ware van de natuur en combineert twee traditionele vormen van voedselverwerving.

2. **Figuratieve betekenis: ambitieuze zoeker**
Een "jachtvisser" kan iemand zijn die voortdurend op zoek is naar kansen, doelen of mogelijkheden. Het gebruik van "jacht" en "visser" symboliseert de gedreven jacht naar succes (de overtuigende impuls) en het geduldige wachten tot iets toehapt, zoals bij vissen.

3. **Letterlijke combinatie van vaardigheden**
In een letterlijke interpretatie kan de term verwijzen naar iemand die 'jaagt in het water' of een combinatie hanteert van traditionele jachttechnieken en visgereedschap tijdens het beoefenen van zijn ambacht. Dit kan wijzen op een innovatieve manier van voedsel zoeken, mogelijk in een survival-setting.

4. **Natuurromantische of poëtische invulling**
De term kan binnen poëzie of literatuur worden ingezet als beeldspraak voor een rusteloze ziel, iemand die balanceert tussen actie (jacht) en geduld (visser), symbolisch voor een mens die zoekt naar balans in een turbulente wereld.

Vis vangt visser met een haak laat hem naar kucht happen en flikkert hem zwaar gewond tussen de wolven op de wal. Dus draai de rollen om wat als jij de vis zou zijn hoe zou jij het vinden
